% ch05-systems.tex -- single-cycle RV32I core in Verilog (capstone part 1)
\chapter{A Single-Cycle RV32I Core}
\label{ch:sc05}

Assembly language is the contract; this chapter builds the machine
that honours it. Everything so far --- the encodings of ch03, the
frames of ch04, the cpu8 of Book 5 --- meets in one Verilog module:
a single-cycle RV32I core that fetches, decodes, executes, touches
memory, and writes back in one clock tick. Book 5 built the blocks
(an 8-bit ALU, a register file, a control decoder); here they are
re-derived at 32 bits against the real ISA --- compute-all-then-mux
ALU, two-read one-write regfile with \texttt{x0} hardwired, and a
decoder that is literally a truth table in Verilog. The core runs a
real program in Icarus: thirteen hand-assembled instructions,
architectural state checked afterward. CPI is 1 by construction ---
the point of the chapter is seeing \emph{why that is a waste of time},
which is the whole argument for the pipeline of ch06.

\section{Fetch--decode--execute, properly}
\label{sec:fde05}
\index{single-cycle processor}\index{datapath}\index{RV32I core}

Ch03 drew the loop as fetch $\to$ decode $\to$ execute $\to$ next.
The honest version splits execute into three: read operands, compute
in the ALU, touch data memory, write the result back --- five acts,
all inside one clock period. The program counter is the only state
element outside the register file and memories: an \texttt{insn} wire
from instruction memory, a \texttt{pc} register, and a next-pc mux.
Everything else is combinational logic evaluating the current
instruction \emph{right now}. That is the definition of
single-cycle: latency of one instruction $=$ one clock period.

\begin{figure}[t]
\centering
\begin{tikzpicture}[font=\scriptsize, >=latex, x=0.85cm, y=0.85cm]
% full single-cycle datapath
\draw[thick, fill=accentsoft] (0,3.4) rectangle (1.6,4.4);
\node[font=\tiny, align=center] at (0.8,3.9) {PC\\reg};
\draw[thick, fill=accentsoft] (2.2,3.4) rectangle (3.8,4.4);
\node[font=\tiny, align=center] at (3.0,3.9) {instr\\mem};
\draw[thick, fill=goodbg] (4.6,3.4) rectangle (6.2,4.4);
\node[font=\tiny, align=center] at (5.4,3.9) {decode\\+ imm};
\draw[thick, fill=goodbg] (7.0,3.4) rectangle (8.6,4.4);
\node[font=\tiny, align=center] at (7.8,3.9) {regfile};
\draw[thick, fill=amberbg] (9.4,3.4) rectangle (11.0,4.4);
\node[font=\tiny, align=center] at (10.2,3.9) {ALU};
\draw[thick, fill=graybg] (11.8,3.4) rectangle (13.4,4.4);
\node[font=\tiny, align=center] at (12.6,3.9) {data\\mem};
\draw[thick, fill=white] (13.9,3.6) circle (0.3);
\node[font=\tiny] at (13.9,3.6) {WB};
\draw[->, very thick] (1.6,3.9) -- (2.2,3.9);
\draw[->, very thick] (3.8,3.9) -- (4.6,3.9);
\draw[->, very thick] (6.2,3.9) -- (7.0,3.9);
\draw[->, very thick] (8.6,3.9) -- (9.4,3.9);
\draw[->, very thick] (11.0,3.9) -- (11.8,3.9);
\draw[->, very thick] (13.4,3.9) -- (13.9,3.9);
% writeback arc
\draw[->, dashed, accent] (13.9,3.3) .. controls (10,2.0) .. (7.8,3.35);
\node[font=\tiny, accent, fill=white, inner sep=1pt] at (10.8,2.45) {wb back to regfile};
% pc+4 loop
\draw[->, dashed, slate] (13.9,3.9) .. controls (12,5.6) and (3,5.6) .. (0.8,4.45);
\node[font=\tiny, slate, fill=white, inner sep=1pt] at (7,5.1) {pc+4 or pc+immB or pc+immJ};
\node[align=left, text=slate, text width=11cm] at (7,0.7)
  {One instruction does the whole journey in one clock: PC $\to$
   fetch $\to$ decode+imm $\to$ read regs $\to$ ALU $\to$ (maybe)
   data mem $\to$ writeback, while the next-PC mux chooses
   pc+4, pc+immB or pc+immJ. Single-cycle is the truth the pipeline
   is trying to fake.};
\end{tikzpicture}

\caption[Single-cycle datapath]{The whole machine. PC feeds
instruction memory; decode splits the word; the register file and
immediate generator feed the ALU; data memory and the writeback mux
finish the journey; the next-PC mux closes the loop.}
\label{fig:datapath05}
\end{figure}

\section{Decode and the immediate generator}
\label{sec:decode05}
\index{immediate generator}\index{decoder}

\begin{example}[decode is a table]\label{ex:ctrl05}
\texttt{ch05\_rv32i\_decode.v} has no cleverness: a \texttt{case}
on the 7-bit opcode, defaults that make an unknown word a NOP. For
\texttt{sw} ($\texttt{op}=0100011$): \texttt{alusrc=1, memwrite=1,
aluop=add} --- base register plus immediate for the address,
\texttt{rd2} already carries the value to store. The truth-table
figure below is the same table drawn.
\end{example}

\begin{figure}[t]
\centering
\setlength{\tabcolsep}{3pt}
\begin{tabularx}{\linewidth}{@{}l c c c c c c X@{}}
\multicolumn{8}{l}{\footnotesize\bfseries Control truth table (the decoder, as a table)}\\[2pt]
\toprule
insn & regwrite & alusrc & memread & memwrite & wbsel & branch/jump & aluop \\
\midrule
R-type (add\ldots) & 1 & 0 & 0 & 0 & alu & -- & f7,f3 \\
I-type (addi\ldots) & 1 & 1 & 0 & 0 & alu & -- & f3 \\
lw & 1 & 1 & 1 & 0 & mem & -- & add \\
sw & 0 & 1 & 0 & 1 & -- & -- & add \\
beq/bne & 0 & 0 & 0 & 0 & -- & branch & sub \\
lui & 1 & 1 & 0 & 0 & immU & -- & -- \\
jal & 1 & -- & 0 & 0 & pc+4 & jump & -- \\
\bottomrule
\multicolumn{8}{p{0.92\linewidth}}{\scriptsize Every row of
\texttt{ch05\_rv32i\_decode.v} is one line of this table --- the
table IS the design.}
\end{tabularx}

\caption[Control truth table]{The decoder as a table. Each row is a
\texttt{case} arm; each column is a control wire. Reading hardware
like this is a transferable skill --- real processors' ``control''
is a much bigger version of the same idea.}
\label{fig:control05}
\end{figure}

\begin{example}[the immediate mux]\label{ex:immgen05}
The five formats of ch03 produce five bit-shuffles: \texttt{immI}
sign-extends \texttt{insn[31:20]}; \texttt{immS} splices
\texttt{[31:25]} and \texttt{[11:7]}; \texttt{immB} and
\texttt{immJ} reorder the same window into jump offsets;
\texttt{immU} just shifts. One mux selects among them on opcode ---
the scrambler is about fifteen lines of wiring.
\end{example}

\begin{figure}[t]
\centering
\begin{tikzpicture}[font=\scriptsize, >=latex]
% immediate generator bit-field diagram
\node[font=\footnotesize\bfseries, align=left] at (5.6,5.2)
  {immediate generator: five shapes from one instruction word};
\foreach \y/\name/\defn in {
  4.4/{I}/{sext insn[31:20]},
  3.6/{S}/{sext insn[31:25]$|$insn[11:7]},
  2.8/{B}/{sign $\times$19 $|$ insn[31] $|$ insn[7] $|$
   insn[30:25] $|$ insn[11:8] $|$ 0},
  2.0/{U}/{insn[31:12]$|$12$\times$0},
  1.2/{J}/{sign $\times$11 $|$ insn[31] $|$
   insn[19:12] $|$ insn[20] $|$ insn[30:21] $|$ 0}}{
  \draw[thick, fill=accentsoft] (0.6,\y-0.28) rectangle (1.6,\y+0.28);
  \node[font=\footnotesize] at (1.1,\y) {\name};
  \node[anchor=west, font=\tiny] at (1.8,\y) {\defn};}
\node[align=left, text=slate, text width=11cm] at (6.2,0.2)
  {Every field borrows bits from fixed positions --- the ISA keeps
   rs1/rs2/rd in place across formats so decode can split off while
   the imm-gen scrambles this wiring in parallel. One mux, five
   inputs, selected by opcode: that is the whole module.};
\end{tikzpicture}

\caption[Immediate generator]{Five immediates from one word. The ISA
keeps operand fields in fixed positions so the register file reads
in parallel with this scrambler --- decode never serialises.}
\label{fig:immgen05}
\end{figure}

\section{Register file and ALU --- Book 5, widened}
\label{sec:reused05}
\index{register file}\index{ALU}

\begin{example}[the cpu8 mapping, honoured]\label{ex:mapping05}
Book 5's ch14 ALU chose its result by mux after computing all ops;
\texttt{ch05\_rv32i\_alu.v} is the same shape at 32 bits, nine ops
(add, sub, and, or, xor, sll, srl, slt, sltu). The regfile keeps
Book 5's
two-read/one-write shape and adds the RV32I rule that made
\texttt{x0} famous: \texttt{(ra1==0)\ ?\ 0\ :\ x[ra1]} \emph{and}
\texttt{we \&\& wa!=0} --- write \texttt{x0} and the write is
silently dropped, which is exactly how the ISA wants it.
\end{example}

\section{Loads, stores and the data-memory interface}
\label{sec:mem05}
\index{load/store}\index{data memory}

Both \texttt{lw} and \texttt{sw} run the ALU as an address adder:
\texttt{rd1} (base) plus the immediate. The difference is which end
of data memory lights up --- \texttt{memread} routes
\texttt{dmem\_rdata} through the writeback mux into \texttt{rd};
\texttt{memwrite} drives \texttt{dmem\_we} with \texttt{rd2} on the
data pins. Addressing mode coverage is deliberately tiny: RV32I has
exactly one addressing mode (base+imm), which is why real ISAs are
easier to teach than they look.

\begin{checkpoint}[the one addressing mode]\label{cp:addr05}
Every RV32I memory instruction is \texttt{offset(base)}. Indexed,
auto-increment and PC-relative modes from bigger ISAs simply don't
exist --- simplicity is the point of RISC. Predict how a C
\texttt{p[i]} loop compiles: \texttt{lw} with a computed base, or
\texttt{addi} the base between iterations. Ch03's arrsum did the
second.
\end{checkpoint}

\section{Running it: 13 words, checked state}
\label{sec:run05}

\begin{example}[the first real program]\label{ex:run05}
\texttt{ch05\_prog.hex} is thirteen words produced by the book's own
assembler --- arithmetic, a store, a load back, a taken branch, a
\texttt{lui}. The testbench plays instruction and data memory around
\texttt{rv32i\_sc\_top}, runs to the terminal \texttt{jal self}, then
inspects architectural state:
\begin{lstlisting}
ran 12 cycles, pc parked at 00000030
  ok: a0 = 1
  ok: a1 = 2
  ok: a2 = a0+a1 = 3
  ok: mem[8] = 3 (sw worked)
  ok: a3 = lw back = 3
  ok: a4 = a3+5 = 8 (beq skipped 99)
  ok: a5 = a4-a2 = 5
  ok: a6 = lui 0x12345000
CPI check: 12 instrs retired in 12 cycles (CPI=1)
ch05 PASS
\end{lstlisting}
Twelve retired in twelve cycles --- the branch skipped one. Real
verilog, real execution, CPI=1 measured, not asserted.
\end{example}

\begin{example}[sim speed as a measurement]\label{ex:simspeed05}
\texttt{ch05\_spin\_tb} runs a 256-iteration arithmetic loop ---
2000 cycles of core. \texttt{time vvp} reports \textbf{9\,ms
wall-clock}: $\approx$220\,k retired instructions per second under
Icarus on this laptop. Simulation speed is why real CPU verification
needs FPGA prototypes --- and why our testbenches stay small and
self-checking.
\end{example}

\section{Why single-cycle wastes time}
\label{sec:critpath05}
\index{critical path}\index{CPI}

\begin{example}[the \texttt{lw} tax]\label{ex:crit05}
\texttt{lw} travels fetch $\to$ decode+regfile $\to$ ALU $\to$ data
memory $\to$ writeback --- every hardware stop in one period. If each
stage costs about the same delay, the clock period must be $5\tau$.
\texttt{add} finishes in $3\tau$ and idles $2\tau$; \texttt{sw}
idles one. The critical-path figure shows the idle tail: single-cycle
CPI=1 is honest, but the \emph{clock} is five times slower than any
stage needs. Pipelining's promise is exactly that reclaimed tail.
\end{example}

\begin{figure}[t]
\centering
\begin{tikzpicture}[font=\scriptsize, >=latex]
% critical path: lw is the longest single-cycle journey
\node[font=\footnotesize\bfseries] at (5.8,4.9) {critical path ---
\texttt{lw} has to do everything};
\draw[thick, fill=accentsoft] (0.5,3.6) rectangle (2.1,4.3);
\node[font=\tiny, align=center] at (1.3,3.95){PC + imem};
\draw[thick, fill=goodbg] (2.9,3.6) rectangle (4.5,4.3);
\node[font=\tiny, align=center] at (3.7,3.95){decode + regfile};
\draw[thick, fill=amberbg] (5.3,3.6) rectangle (6.9,4.3);
\node[font=\tiny, align=center] at (6.1,3.95){ALU (address)};
\draw[thick, fill=graybg] (7.7,3.6) rectangle (9.3,4.3);
\node[font=\tiny, align=center] at (8.5,3.95){data mem read};
\draw[thick, fill=accentsoft] (10.1,3.6) rectangle (11.7,4.3);
\node[font=\tiny, align=center] at (10.9,3.95){WB mux + setup};
\draw[->, very thick] (2.1,3.95)--(2.9,3.95);
\draw[->, very thick] (4.5,3.95)--(4.9,3.95);
\draw[->, very thick] (6.9,3.95)--(7.7,3.95);
\draw[->, very thick] (9.3,3.95)--(10.1,3.95);
\node[align=left, text=slate, text width=11cm] at (6.2,2.6)
  {Five hardware stops in ONE period: fetch, decode+regread, ALU,
   memory, writeback. \texttt{add} skips the memory stop;
   \texttt{sw} skips writeback. The clock must be slow enough for
   the worst case, so every cheaper instruction waits at the end of
   its own cycle --- that idle tail is what pipelining reclaims.};
\end{tikzpicture}

\caption[Critical path]{\texttt{lw} sets the clock period. Every
cheaper instruction finishes early and waits --- the wasted tail is
what five stages of pipeline registers turn into throughput.}
\label{fig:critpath05}
\end{figure}

\begin{figure}[t]
\centering
\begin{tikzpicture}[font=\scriptsize, >=latex]
% sw cycle flow strip
\node[font=\footnotesize\bfseries] at (5.6,4.6) {one \texttt{sw x5,0(x6)} ---
what each block is doing in its single cycle};
\foreach \x/\t in {
  1.0/{fetch insn},
  3.0/{decode: sw $\to$ memwrite=1},
  5.2/{regfile: rd1=x6, rd2=x5},
  7.6/{ALU: addr = x6 + immS},
  10.0/{dmem[addr] $\leftarrow$ x5},
  12.6/{pc $\leftarrow$ pc+4}}{
  \draw[thick, fill=accentsoft] (\x-0.85,3.2) rectangle (\x+0.85,4.0);
  \node[font=\tiny, align=center] at (\x,3.6) {\t};}
\draw[->, very thick] (1.85,3.6)--(2.15,3.6);
\draw[->, very thick] (3.85,3.6)--(4.35,3.6);
\draw[->, very thick] (6.05,3.6)--(6.75,3.6);
\draw[->, very thick] (8.45,3.6)--(9.15,3.6);
\draw[->, very thick] (10.85,3.6)--(11.75,3.6);
\node[align=left, text=slate, text width=11.5cm] at (6.4,2.2)
  {No writeback, no branch: \texttt{sw} leaves regfile-write
   and the writeback mux dark. The control word is the whole
   behaviour --- once the table row is read, the datapath is
   a fixed route the data takes.};
\end{tikzpicture}

\caption[sw cycle flow]{One \texttt{sw} through the datapath: the
control word is read once, then the data takes a fixed route.}
\label{fig:cycleflow05}
\end{figure}

\begin{checkpoint}[predict CPI=1's cost]\label{cp:cpi05}
If the five stages each cost $\tau$ and the register file adds
$\delta$, what's the best achievable clock for the single-cycle core
vs a 5-stage pipeline on the same blocks? $(5\tau+\delta)$ vs
$(\tau+\delta)$: the pipeline's ceiling is the stage time plus
register overhead --- the theme of ch06.
\end{checkpoint}

\section{Common mistakes}
\label{sec:mistakes05}
\begin{itemize}
\item \textbf{Writing \texttt{x0}.} Forgetting \texttt{wa!=0} makes
      \texttt{li x0,5} stick --- every ``hardwired'' invariant breaks.
\item \textbf{Branch on the wrong field.} \texttt{immB} needs the
      shifted splice, not \texttt{immI}: a \texttt{beq} that reads
      \texttt{immI} lands two bytes off silently.
\item \textbf{Muxing before reading.} Using \texttt{rd2} for the
      ALU-B input unconditionally breaks \texttt{lw}/\texttt{sw} ---
      the register file reads the address \emph{register}, the
      immediate is the offset.
\item \textbf{Testing only R-types.} A decode bug in \texttt{sw}
      shows nothing until memory is checked --- the testbench inspects
      \texttt{dmem}, not just registers.
\item \textbf{Believing CPI=1 is fast.} It is honest and it is slow:
      latency hidden by a slow clock is still latency.
\end{itemize}

\section{Summary}
\label{sec:summary05}
\begin{tabularx}{\linewidth}{@{}l X@{}}
\toprule
\textbf{Block} & \textbf{What it is} \\
\midrule
PC + next-pc mux & only state besides regfile/memories \\
decode & opcode \texttt{case} $\to$ control-word table \\
imm generator & five shuffles, one mux, parallel with regfile \\
regfile + ALU & Book 5 blocks, widened, \texttt{x0} hardwired \\
measured & 13-word program, 8 state checks, CPI=1; $\sim$220k instr/s \\
next & five stage registers $\to$ pipeline (ch06) \\
\bottomrule
\end{tabularx}

\section*{Further reading}
\label{sec:further05}
\begin{itemize}
\item Harris \& Harris, \emph{Digital Design and Computer
      Architecture, RISC-V ed.}, ch7 --- the single-cycle
      datapath this chapter mirrors:\\
      \url{https://pages.hmc.edu/harris/ddca/ddcarv.html}
\item Patterson \& Hennessy, \emph{Computer Organization and Design},
      ch4 (the processor):\\
      \href{https://www.elsevier.com/books/computer-organization-and-design-risc-v-edition/patterson/978-0-12-820331-6}{elsevier.com --- COD RISC-V edition}
\item RISC-V unprivileged ISA spec (the contract implemented):\\
      \url{https://github.com/riscv/riscv-isa-manual}
\item CS:APP ch4 (SEQ $\to$ PIPE pedagogical shape):\\
      \url{https://csapp.cs.cmu.edu/}
\item Book 5 ch14--16 (the ALU/regfile/control the code reuses).
\end{itemize}

\begin{rulebox}[Watch instead]
\begin{itemize}
\item \textbf{DDCA Ch7 Part 1: Microarchitecture Introduction} ---
      Sarah Harris, 6:45 --- why the datapath looks the way it does.
      \url{https://youtu.be/lrN-uBKooRY}
\item \textbf{DDCA Ch7 Part 2: RISC-V Single-Cycle Processor
      Datapath: lw} --- Sarah Harris, 7:25 --- the same critical-path
      walk, drawn live. \url{https://youtu.be/AoBkibslRBM}
\end{itemize}
\end{rulebox}

\section*{Notes}
\label{sec:notes05}
\begin{itemize}
\item Datapath/control/imm-gen design follows COD ch4 and DDCA ch7;
      adapted to our rv32i\_* module split.
\item All simulation output captured from Icarus Verilog 13.0 on the
      dev box (arm64 macOS); program bytes assembled by
      \texttt{ch03\_rv32asm.py}.
\item \texttt{bge}/\texttt{blt} not yet decoded --- the program uses
      \texttt{beq}; wider branch support lands with ch07's hazard
      work.
\end{itemize}

\section{Exercises}
\label{sec:exercises05}

\begin{exercise}[easy]\label{ex:e05a}
Which control wires does \texttt{jal} assert, and why is
\texttt{wbsel=pc+4}?
\end{exercise}
\begin{solution}[hint]
regwrite + jump; the destination register must capture the return
address.
\end{solution}

\begin{exercise}[easy]\label{ex:e05b}
Why is \texttt{x0}'s write-enable explicitly masked in the regfile
rather than trusted to software?
\end{exercise}
\begin{solution}[hint]
Hardware invariant: a stray \texttt{sw}-style bug or bad codegen
must not corrupt the constant zero.
\end{solution}

\begin{exercise}[medium]\label{ex:m05a}
Trace \texttt{lw a3,8(zero)} through every block of the datapath:
what does each output wire carry?
\end{exercise}
\begin{solution}[hint]
PC=0x10; decode: alusrc=1, memread=1, wbsel=mem; rd1=0;
ALU=8; dmem[2]=3; wb=3 $\to$ a3.
\end{solution}

\begin{exercise}[medium]\label{ex:m05b}
The imm generator scrambles five shapes in parallel with the regfile
read. Why is the parallelism safe even when an instruction doesn't
need an immediate?
\end{exercise}
\begin{solution}[hint]
The mux only selects when alusrc=1; wasted scrambling costs power,
not correctness.
\end{solution}

\begin{exercise}[medium]\label{ex:m05c}
\texttt{sub} and \texttt{beq} share \texttt{aluop=sub}. What differs
in their control words and why can the ALU be shared?
\end{exercise}
\begin{solution}[hint]
beq sets branch, no regwrite; the subtract output is used only as
the zero flag --- ALU is a function of inputs, free to share.
\end{solution}

\begin{exercise}[medium]\label{ex:m05d}
Our core ran 12 instructions in 12 cycles but the program was 13
words. Reconcile.
\end{exercise}
\begin{solution}[hint]
The taken \texttt{beq} skipped \texttt{addi a4,zero,99} --- fetched
but never decoded... actually never fetched: branch redirect at
execute. The self-loop \texttt{jal} keeps re-executing.
\end{solution}

\begin{exercise}[medium]\label{ex:m05e}
Estimate the same-loop sim speed if the tb also dumped a VCD ---
does the core get slower, or only the log?
\end{exercise}
\begin{solution}[hint]
Both: event-driven sim slows with more dumped signals; the RTL is
unchanged. Measurement vs mechanism.
\end{solution}

\begin{exercise}[hard]\label{ex:h05a}
Add \texttt{slti} and \texttt{slt} to the decoder: which control
fields change, which don't, and how would you test each?
\end{exercise}
\begin{solution}[hint]
\texttt{slti}: I-type row, \texttt{f3=010$\to$aluop=slt}. Test:
\texttt{slti a0,a0,0} should set a0=0 (1$\nless$0) then
\texttt{slti a0,a0,2} = 1.
\end{solution}

\begin{exercise}[hard]\label{ex:h05b}
The critical-path argument assumed equal stage delays. If decode is
actually 2$\tau$, which instruction is now critical and does the
single-cycle case get better or worse relative to pipelining?
\end{exercise}
\begin{solution}[hint]
Still \texttt{lw} (it's the superset path); single-cycle gets
\emph{worse} --- its period grows while the pipeline's is set by the
slowest stage.
\end{solution}

\begin{exercise}[hard]\label{ex:h05c}
Design a self-checking test that proves \texttt{x0} is hardwired:
write to it via \texttt{addi x0,x0,7} then read it. What else must
the tb check to catch a lazy \texttt{we}?
\end{exercise}
\begin{solution}[hint]
x0 must still read 0 after the write; also check a real write to a
nonzero reg landed (proving we wasn't masked globally).
\end{solution}

% chapter end
